% Analytic teaching example: A={1,2,3,4}, B={3,4,5}.
% Three disjoint regions are shown explicitly; no random data or transparency.
\begin{tikzpicture}[x=1mm,y=1mm,font=\figurefont\fontsize{11}{13}\selectfont,
                    text=FigureInk,line width=.75pt]
  \definecolor{OverlapBlue}{HTML}{7188A3}
  \definecolor{OverlapBlueFill}{HTML}{E5EAF0}
  \definecolor{OverlapGreen}{HTML}{829C8B}
  \definecolor{OverlapGreenFill}{HTML}{E8EEE9}
  \definecolor{OverlapYellowFill}{HTML}{F1ECDD}
  \path[use as bounding box] (-46,-29) rectangle (46,31);
  \fill[OverlapBlueFill] (-15,0) circle[radius=25];
  \fill[OverlapGreenFill] (15,0) circle[radius=25];
  \begin{scope}
    \clip (-15,0) circle[radius=25];
    \fill[OverlapYellowFill] (15,0) circle[radius=25];
  \end{scope}
  \draw[OverlapBlue] (-15,0) circle[radius=25];
  \draw[OverlapGreen] (15,0) circle[radius=25];
  \node[text=OverlapBlue] at (-31,26) {$A$};
  \node[text=OverlapGreen] at (31,26) {$B$};
  \node at (-23,7) {$1$};
  \node at (-23,-7) {$2$};
  \node at (0,7) {$3$};
  \node at (0,-7) {$4$};
  \node at (23,0) {$5$};
\end{tikzpicture}
